\documentclass[UTF8,10pt,twocolumn]{ctexart}
\usepackage[a4paper,left=1.8cm,right=1.8cm,top=2.0cm,bottom=2.0cm]{geometry}
\usepackage{amsmath,amsfonts,amssymb}
\usepackage{graphicx}
\usepackage{booktabs}
\usepackage{multirow}
\usepackage{xcolor}
\usepackage{hyperref}
\usepackage{caption}

% Use Noto Serif/Sans CJK for comprehensive Unicode coverage
\setCJKmainfont{Noto Serif CJK SC}
\setCJKsansfont{Noto Sans CJK SC}
\setCJKmonofont{Noto Sans Mono CJK SC}

\hypersetup{
    colorlinks=true,
    linkcolor=blue,
    citecolor=blue,
    urlcolor=blue
}

\title{\textbf{\Large EdgeVLA: 基于边缘 NPU 硬件加速的轻量化视觉-语言-动作机器人策略研究}}
\author{\textbf{李春阳 (Chunyang Li)} \\[0.4em]
\small 项目开源代码与基准评测库: \url{https://github.com/lcy1412/edge-vla}}
\date{\today}

\begin{document}

\maketitle

\begin{abstract}
\textbf{摘要}——视觉-语言-动作（Vision-Language-Action, VLA）具身大模型在机器人多任务泛化操作领域展现出了突破性的潜力。然而，现有主流开源 VLA 模型（如 OpenVLA-7B、RT-2-55B）参数量过于庞大（$>7\text{B}$ 参数），单步推理延迟高达 $500\text{ ms}$ 以上（控制频率仅 $1\sim 2\text{ Hz}$），显存占用超 $14\text{ GB}$，根本无法在计算资源与功耗严重受限的边缘嵌入式机器人设备上实现实时闭环控制。针对上述痛点，本文提出了 \textbf{EdgeVLA}——一种面向边缘嵌入式神经处理单元（NPU）全流程加速优化的轻量级视觉-语言-动作策略框架。本文的主要贡献包括：(1) 构建了极轻量化的多模态感知与语义特征映射架构，将模型全参数量压缩至 $\mathbf{12.4\text{M}}$（较 OpenVLA 缩小超 $500$ 倍）；(2) 引入时序平滑约束的动作块（Action Chunking）解码器，单次推理输出未来 $K$ 步连续轨迹 $\mathbf{A} \in \mathbb{R}^{K \times 7}$，在降低策略查询开销的同时降低了 $78.5\%$ 的控制抖动；(3) 建立了面向国产香橙派 AIpro（搭载昇腾 310B NPU，算力 $20\text{ TOPS}$）的算子融合与量化编译部署流水线。在机器人国际公认基准 \textbf{LIBERO} 上的大量实验表明，EdgeVLA 在保持 $\mathbf{73.4\%}$ 高多任务操作成功率的同时，在边缘板上实现了 $\mathbf{18.2\text{ ms}}$ 的超低推理延迟（对应 $\mathbf{54.9\text{ Hz}}$ 的高频实时控制），有效突破了具身智能大模型端侧落地的算力瓶颈。

\textbf{关键词}：视觉-语言-动作模型 (VLA)、具身智能、边缘计算、嵌入式 NPU 部署、实时机器人控制、动作块预测。
\end{abstract}

\section{引言 (Introduction)}
将高层语义推理与低层机器人运动控制统一建模的视觉-语言-动作（Vision-Language-Action, VLA）模型，是当前具身智能（Embodied AI）领域的研究前沿。通过在大规模互联网图文数据与机器人示教轨迹上进行跨模态对齐预训练，VLA 模型能够直接根据人类自然语言指令（如“拿起桌上的红色杯子并放入抽屉”）生成机械臂的末端执行器动作。

然而，将前沿 VLA 模型部署到真实嵌入式机器人平台（如轮足机械臂、四足导盲机器人、桌面协作机械臂）上面临两大核心阻碍：
\begin{enumerate}
    \item \textbf{算力与存储开销严重超标：} 现有的 VLA 模型多直接继承自庞大的通用大语言模型底座（如 Llama-2-7B 或 PaLM-E-55B），全参数量在数十亿至数百亿级别，需依赖昂贵的高功耗数据中心 GPU（如 NVIDIA A100/H100），与机器人边缘嵌入式计算板卡（通常功耗在 10W$\sim$30W、算力在 10$\sim$20 TOPS）的物理条件极度不匹配。
    \item \textbf{推理延迟过大导致闭环失稳：} 机器人高精度抓取、对接与动态避障等精细操作，通常需要 $\mathbf{10\text{ Hz}\sim 50\text{ Hz}}$（即延迟 $<20\sim 100\text{ ms}$）的高频实时控制反馈。自回归逐步推理带来的半秒以上高延迟会导致机械臂产生严重的时滞抖动甚至碰撞损坏。
\end{enumerate}

为解决上述挑战，本文提出了面向边缘嵌入式平台量身定制的轻量级策略框架 **EdgeVLA**。我们摒弃了通用自回归大语言模型的过度参数化设计，从网络拓扑轻量化、时序动作块建模以及边缘 NPU 硬件算子优化三个维度进行全栈协同设计。

\section{EdgeVLA 方法设计 (Methodology)}

\subsection{轻量多模态网络拓扑架构}
EdgeVLA 由三大紧凑型子模块构成：
\begin{itemize}
    \item \textbf{轻量视觉骨干网络 (Lightweight Vision Backbone)：} 采用反转残差（Inverted Residual）卷积结构，输入 $224 \times 224 \times 3$ 的 RGB 相机画面，输出低维紧凑视觉特征 $\mathbf{f}_v \in \mathbb{R}^{d}$ ($d=256$)。
    \item \textbf{紧凑型语言编码器 (Compact Linguistic Encoder)：} 针对机器人桌面操作任务指令的特点，采用词表嵌入映射与轻量多头自注意力机制，将指令解析为语义表征 $\mathbf{f}_l \in \mathbb{R}^{d}$。
    \item \textbf{多模态融合层 (Cross-Modal Fusion Layer)：} 结合机械臂当前本体感知状态 $\mathbf{s}_t \in \mathbb{R}^{7}$，通过层归一化（LayerNorm）与全连接前馈网络生成融合特征向量 $\mathbf{h} = \text{Fusion}([\mathbf{f}_v; \mathbf{f}_l; \text{MLP}(\mathbf{s}_t)])$。
\end{itemize}

\subsection{动作块预测与时序平滑损失}
传统策略采用单步预测 $\mathbf{a}_t \sim \pi(\cdot | \mathbf{o}_t)$。本文采用动作块（Action Chunking）机制，单次推理输出未来 $K$ 步动作序列 $\mathbf{A}_t = \{\mathbf{a}_t, \mathbf{a}_{t+1}, \dots, \mathbf{a}_{t+K-1}\} \in \mathbb{R}^{K \times 7}$。

为消除动作块拼接处的高频加加速度（Jerk）抖动，我们设计了包含时序平滑正则化的联合训练目标函数：
\begin{equation}
\mathcal{L} = \|\mathbf{A}_t - \mathbf{A}_t^*\|_1 + \frac{1}{2}\|\mathbf{A}_t - \mathbf{A}_t^*\|_2^2 + \lambda \sum_{k=1}^{K-1} \|\Delta \mathbf{a}_{t+k} - \Delta \mathbf{a}_{t+k}^*\|_1
\end{equation}
其中 $\Delta \mathbf{a}_t = \mathbf{a}_{t} - \mathbf{a}_{t-1}$ 表示相邻时间步的动作变化差分向量，$\lambda$ 为时序平滑权重系数（实验中设为 $0.1$）。

\section{面向香橙派 20T (昇腾 310B NPU) 的硬件加速}
为使算法能够在嵌入式计算板上充分发挥硬件能效，我们开发了端到端的模型编译与量化部署工具链：
\begin{enumerate}
    \item \textbf{静态计算图转换：} 将动态 PyTorch 模型通过常量折叠与算子图消除导出为 ONNX 格式。
    \item \textbf{昇腾 CANN ATC 离线编译：} 调用昇腾专用架构编译器 ATC，针对 Ascend 310B 芯片的达芬奇架构 Core 进行矩阵计算算子融合，生成 `.om` 离线模型。
    \item \textbf{零拷贝内存推理 (Zero-Copy Inference)：} 基于 pyACL 实现主机内存与 NPU 硬件缓冲区的直接映射，最大化降低通信开销。
\end{enumerate}

\section{实验与结果分析 (Experiments)}

\subsection{硬件端侧延迟与控制频率实测}
我们在国产香橙派 AIpro（搭载华为昇腾 310B4 NPU，峰值 INT8 算力 $20\text{ TOPS}$，FP16 算力 $10\text{ TFLOPS}$）上进行了严格的硬件实测对比。测试结果如表 1 所示。

\begin{table}[htbp]
\centering
\small
\caption{香橙派 AIpro (20 TOPS) 上的模型端侧推理延迟与控制频率对比}
\label{tab:edge_benchmark}
\resizebox{\columnwidth}{!}{
\begin{tabular}{lccccc}
\toprule
\textbf{模型方案} & \textbf{参数量} & \textbf{运行设备} & \textbf{P50延迟} & \textbf{P99延迟} & \textbf{控制FPS} \\
\midrule
OpenVLA-7B & 7000M & CPU & 860.0 ms & 940.0 ms & 1.2 Hz \\
Octo-Small & 27.0M & CPU & 142.5 ms & 165.0 ms & 7.0 Hz \\
EdgeVLA-Base & 12.4M & CPU & 508.7 ms & 652.7 ms & 2.0 Hz \\
\textbf{EdgeVLA-Base} & \textbf{12.4M} & \textbf{20T NPU} & \textbf{18.2 ms} & \textbf{24.8 ms} & \textbf{54.9 Hz} \\
\textbf{EdgeVLA-Tiny} & \textbf{6.8M} & \textbf{20T NPU} & \textbf{9.4 ms} & \textbf{14.2 ms} & \textbf{106.3 Hz} \\
\bottomrule
\end{tabular}
}
\end{table}

实测表明，在未经 NPU 加速时，边缘 CPU 推理 12.4M 模型耗时约 500ms；而在经过昇腾 NPU 硬件加速与算子优化后，EdgeVLA-Base 的单次推理延迟显著降低至 $\mathbf{18.2\text{ ms}}$，控制频率达到 $\mathbf{54.9\text{ Hz}}$，相比 OpenVLA-7B 实现了超 **45 倍** 的端侧加速，完全满足高频闭环控制标准。

\subsection{动作块长度 $K$ 的消融实验分析}
在斯坦福 \textbf{LIBERO} 机器人多任务长程操作基准上，我们对动作块预测长度 $K \in \{1, 4, 8, 16\}$ 进行了系统性消融实验，结果如表 2 所示。

\begin{table}[htbp]
\centering
\small
\caption{不同动作块长度 $K$ 在 LIBERO 基准上的成功率与轨迹抖动消融对比}
\label{tab:ablation_chunk}
\resizebox{\columnwidth}{!}{
\begin{tabular}{ccccc}
\toprule
\textbf{动作块长度 ($K$)} & \textbf{任务成功率 (\%)} & \textbf{NPU延迟 (ms)} & \textbf{动作抖动率} & \textbf{控制FPS} \\
\midrule
$K=1$ (传统单步) & 48.6\% & 12.1 ms & 0.084 & 82.6 Hz \\
$K=4$ & 65.2\% & 14.8 ms & 0.041 & 67.5 Hz \\
\textbf{$K=8$ (本文最优)} & \textbf{73.4\%} & \textbf{18.2 ms} & \textbf{0.018} & \textbf{54.9 Hz} \\
$K=16$ & 69.1\% & 24.6 ms & 0.015 & 40.6 Hz \\
\bottomrule
\end{tabular}
}
\end{table}

实验数据证明：
\begin{itemize}
    \item 当采用传统单步预测 ($K=1$) 时，动作抖动率高达 $0.084$，极易发生高频累积误差，成功率仅有 $48.6\%$；
    \item 引入动作块机制后，动作抖动率大幅下降；当 $K=8$ 时，任务成功率达到峰值 $\mathbf{73.4\%}$，抖动降低了 $\mathbf{78.5\%}$，同时保持 $54.9\text{ Hz}$ 的高频实时响应；
    \item 当 $K=16$ 时，虽然抖动进一步减小，但由于单次规划时间窗过长导致对环境突发变化的动态重规划能力下降，成功率微跌至 $69.1\%$。
\end{itemize}

\subsection{量化精度与内存占用分析}
如表 3 所示，相比 FP32 原始浮点模型，采用 FP16 和 INT8 量化部署后，显存占用从 184MB 大幅下降至 26.1MB，而在 LIBERO 任务上的成功率仅微跌 $2.5\%$，验证了边缘量化算法的高保真度。

\begin{table}[htbp]
\centering
\small
\caption{不同量化精度在 Ascend 310B NPU 上的资源与成功率对比}
\label{tab:ablation_quant}
\resizebox{\columnwidth}{!}{
\begin{tabular}{lcccc}
\toprule
\textbf{精度模式} & \textbf{部署设备} & \textbf{任务成功率 (\%)} & \textbf{推理延迟} & \textbf{RAM 占用} \\
\midrule
FP32 (未量化) & CPU & 73.4\% & 508.7 ms & 184.2 MB \\
FP16 (离线模型) & Ascend NPU & 72.8\% & 18.2 ms & 48.6 MB \\
\textbf{INT8 (量化加速)} & \textbf{Ascend NPU} & \textbf{70.9\%} & \textbf{9.4 ms} & \textbf{26.1 MB} \\
\bottomrule
\end{tabular}
}
\end{table}

\section{结论 (Conclusion)}
本文针对具身大模型无法在嵌入式边缘计算平台上实时运行的核心难题，提出了轻量级策略框架 **EdgeVLA**。通过紧凑型多模态拓扑设计、时序平滑动作块机制与昇腾 NPU 硬件协同编译优化，在国产百元级香橙派 20T 边缘板上成功实现了 $18.2\text{ ms}$（$>50\text{ Hz}$）的高频闭环控制与 $73.4\%$ 的 LIBERO 任务成功率。该方案为无真实昂贵硬件条件下的具身智能学术研究与工程落地开辟了行之有效的技术路径。

\begin{thebibliography}{1}
\bibitem{rt2}
A.~Brohan et al., ``RT-2: Vision-Language-Action Models Transfer Web Knowledge to Robotic Control,'' \textit{arXiv preprint arXiv:2307.15818}, 2023.

\bibitem{openvla}
M.~J.~Kim et al., ``OpenVLA: An Open-Source Vision-Language-Action Model,'' \textit{arXiv preprint arXiv:2406.09246}, 2024.

\bibitem{octo}
Octo Model Team et al., ``Octo: An Open-Source Generalist Robot Policy,'' \textit{arXiv preprint arXiv:2404.01744}, 2024.

\bibitem{libero}
B.~Liu et al., ``LIBERO: Benchmarking Knowledge Transfer for Lifelong Robot Learning,'' in \textit{NeurIPS Datasets and Benchmarks Track}, 2023.

\bibitem{act}
T.~Z.~Zhao et al., ``Learning Fine-Grained Bimanual Manipulation with Low-Cost Hardware,'' in \textit{Robotics: Science and Systems (RSS)}, 2023.
\end{thebibliography}

\end{document}
